% Options for packages loaded elsewhere
\PassOptionsToPackage{unicode}{hyperref}
\PassOptionsToPackage{hyphens}{url}
\PassOptionsToPackage{dvipsnames,svgnames,x11names}{xcolor}
\documentclass[
  10pt,
  fontset=fandol]{ctexart}
\usepackage{xcolor}
\usepackage[margin=16mm]{geometry}
\usepackage{amsmath,amssymb}
\setcounter{secnumdepth}{-\maxdimen} % remove section numbering
\usepackage{iftex}
\ifPDFTeX
  \usepackage[T1]{fontenc}
  \usepackage[utf8]{inputenc}
  \usepackage{textcomp} % provide euro and other symbols
\else % if luatex or xetex
  \usepackage{unicode-math} % this also loads fontspec
  \defaultfontfeatures{Scale=MatchLowercase}
  \defaultfontfeatures[\rmfamily]{Ligatures=TeX,Scale=1}
\fi
\usepackage{lmodern}
\ifPDFTeX\else
  % xetex/luatex font selection
\fi
% Use upquote if available, for straight quotes in verbatim environments
\IfFileExists{upquote.sty}{\usepackage{upquote}}{}
\IfFileExists{microtype.sty}{% use microtype if available
  \usepackage[]{microtype}
  \UseMicrotypeSet[protrusion]{basicmath} % disable protrusion for tt fonts
}{}
\makeatletter
\@ifundefined{KOMAClassName}{% if non-KOMA class
  \IfFileExists{parskip.sty}{%
    \usepackage{parskip}
  }{% else
    \setlength{\parindent}{0pt}
    \setlength{\parskip}{6pt plus 2pt minus 1pt}}
}{% if KOMA class
  \KOMAoptions{parskip=half}}
\makeatother
\setlength{\emergencystretch}{3em} % prevent overfull lines
\providecommand{\tightlist}{%
  \setlength{\itemsep}{0pt}\setlength{\parskip}{0pt}}
\usepackage{amsmath,amssymb,mathtools,needspace}
\xeCJKDeclareCharClass{CJK}{"2032,"0400->"04FF,"2160->"216B,"2460->"2473,"25A0,"25C7,"25CB,"25CF}
\IfFontExistsTF{Arial Unicode MS}{\xeCJKsetup{AutoFallBack=true}\setCJKfallbackfamilyfont{\CJKrmdefault}{Arial Unicode MS}}{}
\setcounter{secnumdepth}{0}
\setlength{\emergencystretch}{2em}
\allowdisplaybreaks
\usepackage{bookmark}
\IfFileExists{xurl.sty}{\usepackage{xurl}}{} % add URL line breaks if available
\urlstyle{same}
\hypersetup{
  pdftitle={最佳平方逼近},
  colorlinks=true,
  linkcolor={Maroon},
  filecolor={Maroon},
  citecolor={Blue},
  urlcolor={Blue},
  pdfcreator={LaTeX via pandoc}}

\title{最佳平方逼近}
\author{}
\date{}

\begin{document}
\maketitle

\subsection{3.3 最佳平方逼近}\label{ux6700ux4f73ux5e73ux65b9ux903cux8fd1}

用均方误差最小作为度量标准, 研究函数 \(f(x) \in C[a,b]\) 的逼近多项式, 就是这一节要讨论的最佳平方逼近问题. 若存在 \(P_n^*(x) \in H_n\), 使 \[
\|f - P_n^*\|_2 = \sqrt{\int_a^b [f(x) - P_n^*(x)]^2 \mathrm{d}x} = \inf_{P \in H_n} \|f - P\|_2, \tag{3.3.1}
\] \(P_n^*(x)\) 就是 \(f(x)\) 在 \([a,b]\) 上的最佳平方逼近多项式. 这一节将研究 \(P_n^*(x)\) 是否存在, 如何计算 \(P_n^*(x)\), 为此先介绍一些有关内积空间的预备知识.

\subsubsection{3.3.1 内积空间}\label{ux5185ux79efux7a7aux95f4}

定义3.4 设在区间 \((a,b)\) 内, 非负函数 \(\rho(x)\) 满足以下条件, 就称 \(\rho(x)\) 为区间 \((a,b)\) 内的权函数:

\(1^\circ \int_a^b |x|^n \rho(x) \mathrm{d}x (n = 0,1,\cdots)\) 存在;

\(2^\circ\) 对于非负的连续函数 \(g(x)\), 若 \[
\int_a^b g(x) \rho(x) \mathrm{d}x = 0, \tag{3.3.2}
\] 则在 \((a,b)\) 内 \(g(x) \equiv 0\).

定义3.5 设 \(f(x), g(x) \in C[a,b], \rho(x)\) 是 \([a,b]\) 上的权函数, 积分 \[
(f,g) = \int_a^b \rho(x) f(x) g(x) \mathrm{d}x \tag{3.3.3}
\] 称为函数 \(f(x)\) 与 \(g(x)\) 在 \([a,b]\) 上的内积.

容易验证这样定义的内积满足下列四条公理:

\begin{enumerate}
\def\labelenumi{(\arabic{enumi})}
\item
  \((f,g) = (g,f)\);
\item
  \((cf,g) = c(f,g), c\) 为常数;
\item
  \((f_1 + f_2, g) = (f_1, g) + (f_2, g)\);
\item
  \((f,f) \geqslant 0\), 当且仅当 \(f=0\) 时 \((f,f)=0\).
\end{enumerate}

满足内积定义的函数空间称为内积空间, 因此, 连续函数空间 \(C[a,b]\) 上定义了内积就形成一个内积空间. 这里定义的内积是 \(n\) 维 Euclid 空间 \(\mathbf{R}^n\) 中两个向量内积定义的推广, 设 \[
\boldsymbol{f} = (f_1, f_2, \cdots, f_n)^{\mathrm{T}}, \quad \boldsymbol{g} = (g_1, g_2, \cdots, g_n)^{\mathrm{T}},
\]

\href{../../source-images/pdf-066.jpeg}{查看原页}

其内积定义为 \((\boldsymbol{f}, \boldsymbol{g}) = \sum_{k=1}^{n} f_{k} g_{k}\); 向量 \(\boldsymbol{f} \in \mathbf{R}^{n}\) 的模(范数)定义为 \[
\| \boldsymbol{f} \|_{2} = \left( \sum_{k=1}^{n} f_{k}^{2} \right)^{\frac{1}{2}},
\] 将它推广到任何内积空间中就有下面的定义.

定义 3.6 \(f(x) \in C[a, b]\), 量 \[
\| f \|_{2} = \sqrt{\int_{a}^{b} \rho(x) f^{2}(x) \mathrm{d}x} = \sqrt{(f, f)} \tag{3.3.4}
\]

称为 \(f(x)\) 的 Euclid 范数.

它同样满足范数三条性质(见 3.1 节), 头两条是显然的, 下面定理包含了第三条性质和其他重要结论.

定理 3.5 对于任何 \(f, g \in C[a, b]\), 下列结论成立: \(1^{\circ}\)

\[
|(f,g)|\leqslant\|f\|_2\|g\|_2\qquad\text{(Cauchy-Schwarz 不等式)};\tag{3.3.5}
\] \(2^{\circ}\)

\[
\|f+g\|_2\leqslant\|f\|_2+\|g\|_2\qquad\text{(三角不等式)};\tag{3.3.6}
\] \(3^{\circ} \| f + g \|_{2}^{2} + \| f - g \|_{2}^{2} = 2(\| f \|_{2}^{2} + \| g \|_{2}^{2})\) (平行四边形定律).

证明 若 \(g=0\), 则式(3.3.5)显然成立, 现考虑 \(g \neq 0\). 对于任何实数 \(\lambda\), 有 \[
0 \leqslant (f + \lambda g, f + \lambda g) = (f, f) + 2\lambda(f, g) + \lambda^{2}(g, g).
\]

现取 \(\lambda = -\frac{(f, g)}{\| g \|_{2}^{2}}\), 代入上式得 \[
\| f \|_{2}^{2} - 2 \frac{|(f, g)|^{2}}{\| g \|_{2}^{2}} + \frac{|(f, g)|^{2}}{\| g \|_{2}^{2}} \geqslant 0, \quad |(f, g)|^{2} \leqslant \| f \|_{2}^{2} \| g \|_{2}^{2}.
\] 两边开平方即得式(3.3.5).

利用式(3.3.5)考虑 \[
\| f + g \|_{2}^{2} = (f + g, f + g) = (f, f) + 2(f, g) + (g, g)
\] \[
\leqslant \| f \|_{2}^{2} + 2 |(f, g)| + \| g \|_{2}^{2}
\] \[
\leqslant \| f \|_{2}^{2} + 2 \| f \|_{2} \| g \|_{2} + \| g \|_{2}^{2}
\] \[
= (\| f \|_{2} + \| g \|_{2})^{2},
\] 两边开方则得式(3.3.6).

平行四边形定律可直接计算得出, 可作为习题由读者完成. 证毕.

在 \(n\) 维空间中两个向量正交的定义也可推广到内积空间中.

定义 3.7 若 \(f(x), g(x) \in C[a, b]\) 满足 \[
(f, g) = \int_{a}^{b} \rho(x) f(x) g(x) \mathrm{d}x = 0, \tag{3.3.7}
\] 则称 \(f\) 与 \(g\) 在 \([a, b]\) 上带权 \(\rho(x)\) 正交. 若函数族 \(\varphi_{0}(x), \varphi_{1}(x), \cdots, \varphi_{n}(x), \cdots\) 满足关系 \[
(\varphi_{j}, \varphi_{k}) = \int_{a}^{b} \rho(x) \varphi_{j}(x) \varphi_{k}(x) \mathrm{d}x = \begin{cases} 0, & j \neq k, \\ A_{k} > 0, & j = k, \end{cases} \tag{3.3.8}
\]

\href{../../source-images/pdf-067.jpeg}{查看原页}

则称 \(\{\varphi_k\}\) 是 \([a,b]\) 上带权 \(\rho(x)\) 的正交函数族; 若 \(A_k \equiv 1\), 则称 \(\{\varphi_k\}\) 为标准正交函数族. 例如, 三角函数族 \[
1, \cos x, \sin x, \cos 2x, \sin 2x, \cdots
\] 就是在区间 \([-\pi, \pi]\) 上的正交函数族 (权 \(\rho(x) \equiv 1\)), 其 \[
(1,1) = 2\pi,
\] \[
(\sin nx, \sin mx) = \int_{-\pi}^{\pi} \sin nx \sin mx \mathrm{d}x = \begin{cases} \pi, & m=n, \\ 0, & m \neq n \end{cases} (n,m=1,2, \cdots),
\] \[
(\cos nx, \cos mx) = \int_{-\pi}^{\pi} \cos nx \cos mx \mathrm{d}x = \begin{cases} \pi, & m=n, \\ 0, & m \neq n \end{cases} (n,m=1,2, \cdots),
\] \[
(\cos nx, \sin mx) = \int_{-\pi}^{\pi} \cos nx \sin mx \mathrm{d}x = 0 (n,m=1,2, \cdots).
\] 在 \(\mathbf{R}^n\) 空间中, 任一向量都可用它的一组线性无关的基表示, 内积空间的任一元素 \(f(x) \in C[a,b]\) 也同样可用线性无关的基表示, 此时相应地有如下定义.

定义 3.8 设 \(\varphi_0(x), \varphi_1(x), \cdots, \varphi_{n-1}(x)\) 在 \([a,b]\) 上连续, 如果 \[
a_0 \varphi_0(x) + a_1 \varphi_1(x) + \cdots + a_{n-1} \varphi_{n-1}(x) = 0
\] 当且仅当 \(a_0 = a_1 = \cdots = a_{n-1} = 0\) 时成立, 则称 \(\varphi_0, \varphi_1, \cdots, \varphi_{n-1}\) 在 \([a,b]\) 上是线性无关的.

若函数族 \(\{\varphi_k\} (k=0,1, \cdots)\) 中的任何有限个 \(\varphi_k\) 线性无关, 则称 \(\{\varphi_k\}\) 为线性无关函数族.

例如, \(1,\allowbreak  x,\allowbreak  x^2,\allowbreak  \cdots,\allowbreak  x^n,\allowbreak  \cdots\) 就是 \([a,b]\) 上的线性无关函数族. 若 \(\varphi_0(x), \varphi_1(x), \cdots, \varphi_{n-1}(x)\) 是 \([a,b]\) 上的线性无关函数, 且 \(a_0, a_1, \cdots, a_{n-1}\) 是任意实数, 则 \[
S(x) = a_0 \varphi_0(x) + a_1 \varphi_1(x) + \cdots + a_{n-1} \varphi_{n-1}(x)
\] 的全体是 \(C[a,b]\) 中的一个子集, 记作 \[
\Phi = \operatorname{span}\{\varphi_0, \varphi_1, \cdots, \varphi_{n-1}\}. \tag{3.3.9}
\] 下面给出判断函数族 \(\{\varphi_k\} (k=0,1, \cdots, n-1)\) 线性无关的充要条件.

定理 3.6 \(\varphi_0(x), \varphi_1(x), \cdots, \varphi_{n-1}(x)\) 在 \([a,b]\) 上线性无关的充分必要条件是它的 Cramer 行列式 \(G_{n-1} \neq 0\), 其中 \[
G_{n-1} = G(\varphi_0, \varphi_1, \cdots, \varphi_{n-1}) = \left| \begin{array}{cccc} (\varphi_0, \varphi_0) & (\varphi_0, \varphi_1) & \cdots & (\varphi_0, \varphi_{n-1}) \\ (\varphi_1, \varphi_0) & (\varphi_1, \varphi_1) & \cdots & (\varphi_1, \varphi_{n-1}) \\ \vdots & \vdots & & \vdots \\ (\varphi_{n-1}, \varphi_0) & (\varphi_{n-1}, \varphi_1) & \cdots & (\varphi_{n-1}, \varphi_{n-1}) \end{array} \right|. \tag{3.3.10}
\] 定理的证明由读者自行完成.

\subsubsection{3.3.2 函数的最佳平方逼近}\label{ux51fdux6570ux7684ux6700ux4f73ux5e73ux65b9ux903cux8fd1}

下面研究在区间 \([a,b]\) 上一般的最佳平方逼近问题. 对 \(f(x) \in C[a,b]\) 及 \(C[a,b]\)

\begin{quote}
校注（agent 补充，原书术语疑误）：原页定理 3.6 写作``Cramer 行列式''，按原文保留。式 (3.3.10) 是由内积组成的 Gram 行列式，通常称``Gram（格拉姆）行列式''。
\end{quote}

\href{../../source-images/pdf-068.jpeg}{查看原页}

中的一个子集 \(\Phi = \text{span}\{\varphi_0, \varphi_1, \dots, \varphi_n\}\), 若存在 \(S^*(x) \in \Phi\), 使 \[
\| f - S^* \|_2^2 = \inf_{S \in \varphi} \| f - S \|_2^2 = \inf_{S \in \varphi} \int_a^b \rho(x) [f(x) - S(x)]^2 \mathrm{d}x, \tag{3.3.11}
\] 则称 \(S^*(x)\) 是 \(f(x)\) 在子集 \(\Phi \subseteq C[a, b]\) 中的最佳平方逼近函数. 为了求 \(S^*(x)\), 由式 (3.3.11)可知, 该问题等价于求多元函数 \[
I(a_0, a_1, \dots, a_n) = \int_a^b \rho(x) \left[ \sum_{j=0}^n a_j \varphi_j(x) - f(x) \right]^2 \mathrm{d}x \tag{3.3.12}
\] 的最小值. 由于 \(I(a_0, a_1, \dots, a_n)\) 是关于 \(a_0, a_1, \dots, a_n\) 的二次函数, 利用多元函数极值的必要条件 \[
\frac{\partial I}{\partial a_k} = 0 \quad (k = 0, 1, \dots, n),
\] \[
\frac{\partial I}{\partial a_k} = 2 \int_a^b \rho(x) \left[ \sum_{j=0}^n a_j \varphi_j(x) - f(x) \right] \varphi_k(x) \mathrm{d}x = 0 \quad (k = 0, 1, \dots, n),
\] 于是有 \[
\sum_{j=0}^n (\varphi_k, \varphi_j) a_j = (f, \varphi_k) \quad (k = 0, 1, \dots, n). \tag{3.3.13}
\] 这是关于 \(a_0, a_1, \dots, a_n\) 的线性方程组, 称为法方程. 由于 \(\varphi_0, \varphi_1, \dots, \varphi_n\) 线性无关, 故系数行列式 \(G(\varphi_0, \varphi_1, \dots, \varphi_n) \neq 0\), 于是方程组 (3.3.13) 有唯一解 \(a_k = a_k^* (k = 0, 1, \dots, n)\), 从而得到 \[
S^*(x) = a_0^* \varphi_0(x) + \dots + a_n^* \varphi_n(x).
\] 下面证明 \(S^*(x)\) 满足式 (3.3.11), 即对任何 \(S(x) \in \Phi\), 有 \[
\int_a^b \rho(x) [f(x) - S^*(x)]^2 \mathrm{d}x \leqslant \int_a^b \rho(x) [f(x) - S(x)]^2 \mathrm{d}x. \tag{3.3.14}
\] 为此只要考虑 \[
D = \int_a^b \rho(x) [f(x) - S(x)]^2 \mathrm{d}x - \int_a^b \rho(x) [f(x) - S^*(x)]^2 \mathrm{d}x
\] \[
= \int_a^b \rho(x) [S(x) - S^*(x)]^2 \mathrm{d}x + 2 \int_a^b \rho(x) [S^*(x) - S(x)] [f(x) - S^*(x)] \mathrm{d}x.
\] 由于 \(S^*(x)\) 的系数 \(a_k^*\) 是方程 (3.3.13) 的解, 故 \[
\int_a^b \rho(x) [f(x) - S^*(x)] \varphi_k(x) \mathrm{d}x = 0 \quad (k = 0, 1, \dots, n),
\] 从而上式第二个积分为零, 于是 \[
D = \int_a^b \rho(x) [S(x) - S^*(x)]^2 \mathrm{d}x \geqslant 0,
\] 故式 (3.3.14) 成立. 这就证明了 \(S^*(x)\) 是 \(f(x)\) 在 \(\Phi\) 中的最佳平方逼近函数. 若令 \(\delta = f(x) - S^*(x)\), 则平方误差为 \[
\|\delta\|_2^2 = (f - S^*, f - S^*) = (f, f) - (S^*, f) = \|f\|_2^2 - \sum_{k=0}^n a_k^* (\varphi_k, f). \tag{3.3.15}
\]

\begin{quote}
校注（agent 补充，原书符号不一致）：式 (3.3.11) 两处下确界下标原印 \(S\in\varphi\)（小写），本页前文子集定义为 \(\Phi\)（大写）。此处按原图保留小写，所指应为前文的 \(\Phi\)。
\end{quote}

\href{../../source-images/pdf-069.jpeg}{查看原页}

取 \(\varphi_k(x) = x^k, \rho(x) \equiv 1, f(x) \in C[0,1]\), 即要在 \(H_n\) 中求 \(n\) 次最佳平方逼近多项式

\[
S^*(x) = a_0^* + a_1^* x + \cdots + a_n^* x^n,
\]

此时

\[
(\varphi_j, \varphi_k) = \int_0^1 x^{k+j} \mathrm{d}x = \frac{1}{k+j+1},
\]

\[
(f, \varphi_k) = \int_0^1 f(x) x^k \mathrm{d}x \equiv d_k.
\]

若用 \(\boldsymbol{H}\) 表示行列式 \(G_n = G(1, x, x^2, \cdots, x^n)\) 对应的矩阵，则

\[\boldsymbol{H} = \begin{pmatrix}
1 & 1/2 & \cdots & 1/(n+1) \\
1/2 & 1/3 & \cdots & 1/(n+2) \\
\vdots & \vdots & & \vdots \\
1/(n+1) & 1/(n+2) & \cdots & 1/(2n+1)
\end{pmatrix}, \tag{3.3.16}
\]

\(\boldsymbol{H}\) 称为 Hilbert 矩阵. 记

\[
\boldsymbol{a} = (a_0, a_1, \cdots, a_n)^{\mathrm{T}}, \quad \boldsymbol{d} = (d_0, d_1, \cdots, d_n)^{\mathrm{T}},
\]

其中

\[
d_k = (f, x^k) (k=0, 1, \cdots, n), \tag{3.3.17}
\]

则方程

\[\boldsymbol{H}\boldsymbol{a}=\boldsymbol{d}
\]

的解 \(a_k = a_k^*\) (\(k=0, 1, \cdots, n\)) 即为所求.

例 3.2 设 \(f(x) = \sqrt{1+x^2}\), 求 \([0,1]\) 上的一次最佳平方逼近多项式.

解 利用式 (3.3.17), 得

\[
d_0 = \int_0^1 \sqrt{1+x^2} \mathrm{d}x = \frac{1}{2} \ln(1+\sqrt{2}) + \frac{\sqrt{2}}{2} \approx 1.147,
\]

\[
d_1 = \int_0^1 x \sqrt{1+x^2} \mathrm{d}x = \frac{1}{3} (1+x^2)^{\frac{3}{2}} \Bigg|_0^1 = \frac{2\sqrt{2}-1}{3} \approx 0.609,
\]

由方程组

\[
\begin{bmatrix}
1 & \frac{1}{2} \\
\frac{1}{2} & \frac{1}{3}
\end{bmatrix} \begin{pmatrix} a_0 \\ a_1 \end{pmatrix} = \begin{pmatrix} 1.147 \\ 0.609 \end{pmatrix},
\]

解出

\[
a_0 = 0.934, \quad a_1 = 0.426, \quad S_1^*(x) = 0.934 + 0.426x.
\]

平方误差

\[
\|\delta\|_2^2 = (f, f) - (S_1^*, f)
\]

\[
= \int_0^1 (1+x^2) \mathrm{d}x - 0.426 d_1 - 0.934 d_0 = 0.0026,
\]

最大误差

\[
\|\delta\|_\infty = \max_{0 \leqslant x \leqslant 1} |\sqrt{1+x^2} - S_1^*(x)| \approx 0.066.
\]

用 \(\{1, x, x^2, \cdots, x^n\}\) 作基求最佳平方逼近多项式, 当 \(n\) 较大时, 系数矩阵式 (3.3.16) 是高度病态的 (病态矩阵概念见第 7 章), 求法方程 (3.3.13) 的解时, 舍入误差很大, 这时, 用正交多项式作基才能求得最小平方逼近多项式 (见 3.5 节).

\begin{quote}
校注（agent 补充，数值舍入）：原书例 3.2 的系数与平方误差 \(0.0026\) 均按原文保留。该数值来自把 \(d_0,d_1\) 分别近似为 \(1.147,0.609\) 后代入简化式；对于印出的函数 \(S_1^*(x)=0.934+0.426x\)，直接计算 \(\int_0^1[\sqrt{1+x^2}-S_1^*(x)]^2\,\mathrm{d}x\approx0.0007136324\)，与 \(0.0026\) 不同。舍入后的系数不再严格满足精确内积对应的法方程，不能把简化式的数值视为实际平方误差。
\end{quote}

\subsection{本节覆盖原页的校注记录}\label{ux672cux8282ux8986ux76d6ux539fux9875ux7684ux6821ux6ce8ux8bb0ux5f55}

下列记录按原页关联，含同页相邻小节的校注；计算时同时读取上文校注和完整原页。

\begin{itemize}
\tightlist
\item
  \texttt{ch03a-E05}：原书 PDF 页 67
\item
  \texttt{ch03a-E06}：原书 PDF 页 68
\item
  \texttt{ch03a-E07}：原书 PDF 页 69
\end{itemize}

\end{document}
